\chapter{Metodologia}
\label{chap:metodologia}

Este capítulo apresenta a metodologia empírica adotada para investigar a
existência e as propriedades do prêmio de risco de volatilidade no mercado de
Bitcoin, bem como seu conteúdo informacional para a previsão de retornos futuros
do ativo. Com base nas evidências descritivas apresentadas no
Capítulo~\ref{chap:dados}, o capítulo define as questões de pesquisa, as
variáveis, as amostras e os modelos utilizados ao longo da dissertação.

A estratégia empírica tem duas etapas. Na primeira, a dissertação prevê o BVRP
fora da amostra, apenas com informação disponível em $t$. Na segunda, testa se
o BVRP previsto prevê os retornos futuros do Bitcoin, primeiro com um modelo
linear e depois com um modelo não linear em árvore, capaz de capturar
não linearidades, efeitos de regime e interações características do mercado de
criptoativos. Este capítulo concentra-se na definição formal das variáveis, das
amostras e das especificações empíricas; os capítulos seguintes apresentam os
resultados.

% =========================================================
\section{Questões de pesquisa}
\label{sec:met-questoes}

Quatro questões guiam a análise empírica, cada uma tratada em um capítulo:

\begin{enumerate}
    \item \textbf{Existência do prêmio (H1)}: o BVRP é negativo em média? O
          Capítulo~\ref{chap:dados} (Seção~\ref{sec:dados-h1}) responde com um teste
          $t$ sobre a média, com erro-padrão de Newey--West de $h+1$
          defasagens (Seção~\ref{sec:met-inferencia}).
    \item \textbf{Previsão do prêmio}: a informação disponível em $t$ prevê o
          BVRP fora da amostra melhor que a média histórica? O
          Capítulo~\ref{chap:previsao_bvrp} compara modelos lineares e em árvore.
    \item \textbf{Previsibilidade linear dos retornos}: o BVRP previsto prevê
          linearmente os retornos futuros do Bitcoin? O
          Capítulo~\ref{chap:retornos_linear} estima regressões lineares em seis
          horizontes.
    \item \textbf{Previsibilidade não linear}: um modelo em árvore, que escolhe
          sozinho os cortes nas volatilidades, encontra a relação que o modelo
          linear não encontra? O Capítulo~\ref{chap:retornos_nao_linear} responde.
\end{enumerate}

As estratégias de negociação baseadas no prêmio ficam em apêndice
(Apêndice~\ref{app:estrategias}), como análise exploratória.

% =========================================================
\section{Definição das variáveis e horizontes de previsão}
\label{sec:met-variaveis}

Esta seção define formalmente as variáveis utilizadas nos modelos empíricos,
bem como os horizontes temporais considerados na análise. Todas as variáveis
são construídas a partir das séries descritas no Capítulo~\ref{chap:dados}.

A variável dependente dos Capítulos~\ref{chap:retornos_linear}
e~\ref{chap:retornos_nao_linear} é o retorno futuro do Bitcoin no mercado à vista, definido
como o retorno simples acumulado ao longo de um horizonte de $h$ dias:
\[
R_{t+h} = \frac{P_{t+h}}{P_t} - 1,
\]
em que $P_t$ denota o preço de fechamento do Bitcoin no instante $t$ e
$h \in \{1, 5, 10, 20, 30, 60\}$ representa o horizonte de previsão. O retorno
simples corresponde ao retorno efetivo de uma posição no ativo; a volatilidade
histórica, por sua vez, é construída a partir de log-retornos diários,
convenção padrão para agregação temporal de variância.

A variável-alvo do Capítulo~\ref{chap:previsao_bvrp} é o BVRP,
$\text{BVRP}_t = \text{VH}_{t+1:t+30} - \text{IV}_t$, definido na
Seção~\ref{sec:refer-vrp-teorico}. Ele só é conhecido em $t+30$ e, por isso, nunca
entra como variável explicativa. A versão conhecida em $t$ é a proxy,
\[
\text{BVRP}^{\text{proxy}}_t = \text{VH}_{t-29:t} - \text{IV}_t,
\]
que supõe um passeio aleatório sem deriva para a volatilidade histórica: a
melhor previsão da VH dos próximos 30 dias seria a VH dos últimos 30. A
Seção~\ref{sec:refer-vrp-teorico} discute essa hipótese, sabidamente uma
simplificação. Na amostra completa, a correlação entre a proxy e o BVRP na
mesma data é de $-0{,}02$: a proxy aproxima mal a realização do prêmio, o que
motiva prevê-lo com um conjunto mais amplo de variáveis.

As variáveis explicativas, todas conhecidas em $t$, são 17, em quatro grupos:
(i) a volatilidade histórica de 1, 5, 30, 60 e 90 dias, em que as três
primeiras seguem a estrutura diária, semanal e mensal do modelo HAR de
\citet{corsi2009har}; (ii) a volatilidade implícita de 30 dias, suas variações
em 1 e 5 dias e seus desvios em relação às médias móveis de 5 e 30 dias;
(iii) o log-retorno diário, os log-retornos acumulados de 5 e 30 dias e o log
do preço relativo à sua média móvel de 30 dias; e (iv) a proxy, sua variação
diária e o prêmio realizado defasado,
$\text{VH}_{t-29:t} - \text{IV}_{t-30}$, isto é, o BVRP de $t-30$, que se
realiza em $t$. As séries classificadas como I(1) --- o preço e as médias
móveis da IV, em que o ADF não rejeita a raiz unitária e o KPSS rejeita a
estacionariedade --- entram transformadas; as volatilidades, persistentes,
entram em nível. Como $\text{BVRP}^{\text{proxy}}_t = \text{VH}_{t-29:t} -
\text{IV}_t$ por construção, os modelos lineares excluem a VH de 30 dias para
evitar colinearidade perfeita. O Capítulo~\ref{chap:previsao_bvrp} traz o dicionário
completo das variáveis.

% =========================================================
\section{Amostras e janela de estimação}
\label{sec:met-amostras}

A dissertação usa duas amostras. A amostra completa, com 1.806 observações
diárias de 24/03/2021 a 03/03/2026, serve à análise descritiva e ao teste de H1
(Capítulo~\ref{chap:dados}). A amostra de modelagem, com 1.776 observações de
23/04/2021 a 03/03/2026, serve aos Capítulos~\ref{chap:previsao_bvrp},
\ref{chap:retornos_linear} e~\ref{chap:retornos_nao_linear}. Ela começa 30 dias depois porque
o prêmio realizado defasado exige a IV de $t-30$, e o índice DVOL só existe a
partir de 24/03/2021. As duas amostras terminam em 03/03/2026, o último dia
com retorno de 60 dias à frente. O alvo do Capítulo~\ref{chap:previsao_bvrp} só
precisa de 30 dias à frente, mas uma única amostra de modelagem torna os
resultados dos três capítulos comparáveis.

A dissertação estima os modelos de previsão numa subamostra e os avalia fora
dela, em origens de previsão sucessivas. A primeira janela de estimação tem 730 dias,
de 23/04/2021 a 22/04/2023. Em cada origem $t$, o modelo usa apenas
as datas $s \le t - h$ --- o embargo: o alvo de $s$ olha $h$ dias à frente e só
é conhecido em $s + h$ --- e prevê as datas de $t$ até a véspera da origem
seguinte, sempre com as variáveis explicativas da própria data. No
Capítulo~\ref{chap:previsao_bvrp}, $h = 30$, o horizonte do BVRP. O modelo é
reestimado a cada 30 dias. A primeira origem é 21/06/2023, e o período fora da
amostra tem 987 datas (de 21/06/2023 a 03/03/2026), em 33 origens. A primeira
origem fica 60 dias depois do fim da primeira janela, o maior horizonte dos
Capítulos~\ref{chap:retornos_linear} e~\ref{chap:retornos_nao_linear}: com
qualquer $h \le 60$, o treino da primeira origem contém a janela inteira, e o
período fora da amostra é o mesmo em todos os horizontes. Com $h = 30$, o
primeiro treino vai até $t - 30$ e tem 760 observações: as 730 da primeira
janela e as 30 seguintes, cujos alvos já se realizaram. Esse período
é o mesmo nos Capítulos~\ref{chap:previsao_bvrp} a~\ref{chap:retornos_nao_linear}, porque o
regressor dos Capítulos~\ref{chap:retornos_linear} e~\ref{chap:retornos_nao_linear} é o BVRP
previsto no Capítulo~\ref{chap:previsao_bvrp}.

Na janela principal, expansiva, o treino acumula todas as datas desde o início
da amostra de modelagem; como robustez, a janela móvel mantém os últimos 730
dias. Em cada reestimação, a validação cruzada embargada escolhe os
hiperparâmetros dos modelos regularizados e em árvore: cinco dobras expansivas
dentro do treino da origem, com o mesmo embargo --- numa dobra cuja validação
começa em $v$, o treino da dobra usa $s \le v - h$.

% =========================================================
\section{Modelos econométricos}
\label{sec:met-modelos}

Um modelo é uma especificação com suas premissas; o método de mínimos
quadrados ordinários (MQO) é um estimador. Nas especificações lineares desta
seção, o MQO é o estimador padrão, adequado sob a condição de ortogonalidade
$E(\varepsilon \mid X) = 0$, em que $\varepsilon$ é o erro e $X$, o vetor de
regressores.

\subsection{Previsão do prêmio}

O Capítulo~\ref{chap:previsao_bvrp} modela o BVRP como função das variáveis
explicativas,
\begin{equation}
\text{BVRP}_t = f(X_t) + u_t,
\label{eq:met-previsao-bvrp}
\end{equation}
em que $X_t$ reúne as 17 variáveis da Seção~\ref{sec:met-variaveis}. No
modelo linear, $f(X_t) = X_t'\gamma$, estimado por MQO, Ridge ou LASSO; nos
modelos em árvore, $f$ é estimada por floresta aleatória ou por XGBoost
(Capítulo~\ref{chap:referencial}). As referências são a média histórica, a
proxy e o prêmio realizado defasado (Seção~\ref{sec:met-avaliacao}).

\subsection{Retornos e BVRP previsto: modelo linear}

O Capítulo~\ref{chap:retornos_linear} estima, para cada horizonte $h$, o modelo
linear
\begin{equation}
R_{t+h} = \alpha_h + \beta_h \, \widehat{\text{BVRP}}_t + \varepsilon_{t+h},
\label{eq:met-retorno-linear}
\end{equation}
em que $\widehat{\text{BVRP}}_t$ é o BVRP previsto para a data $t$ pelo modelo
estimado na origem vigente, $\alpha_h$ é um intercepto específico ao horizonte
de previsão, $\beta_h$ é o coeficiente de interesse e $\varepsilon_{t+h}$
representa o termo de erro. O estimador é o MQO, aplicado às 987 datas fora da
amostra. A previsão principal vem do Ridge, e a da floresta aleatória serve de
robustez --- escolha fixada antes de ver os resultados. A análise não impõe
restrições \textit{a priori} sobre o sinal de $\beta_h$.

Duas extensões complementam o modelo~\eqref{eq:met-retorno-linear}. A primeira
substitui o prêmio por seus componentes:
\begin{equation}
R_{t+h} = \alpha_h + \beta_{1,h} \, \text{VH}_{t-29:t} + \beta_{2,h} \,
\text{IV}_t + \varepsilon_{t+h}.
\label{eq:met-retorno-componentes}
\end{equation}
Se $\beta_{1,h} + \beta_{2,h} = 0$, a diferença $\text{VH} - \text{IV}$
sintetiza o conteúdo informacional conjunto das duas medidas. A segunda
acrescenta o regime de volatilidade:
\begin{equation}
R_{t+h} = \alpha_h + \beta_h \, \widehat{\text{BVRP}}_t + \delta_h \, D_t
+ \gamma_h \, D_t \, \widehat{\text{BVRP}}_t + \varepsilon_{t+h},
\label{eq:met-retorno-regime}
\end{equation}
em que $D_t = 1$ no regime de alta volatilidade
(Seção~\ref{sec:met-regimes}).

\subsection{Retornos e BVRP previsto: modelo em árvore}

O Capítulo~\ref{chap:retornos_nao_linear} reestima o modelo~\eqref{eq:met-retorno-linear} com
$\widehat{\text{BVRP}}_t = \hat{f}(X_t) - \text{IV}_t$, em que $\hat{f}$ é uma
floresta aleatória que prevê $\text{VH}_{t+1:t+30}$ com as mesmas variáveis,
origens e embargo. A floresta escolhe sozinha os cortes nas volatilidades, o
que substitui os regimes construídos à mão por um único modelo em árvore. Como
robustez, a floresta prevê o BVRP diretamente, e o Ridge entra nas duas
formas.

% =========================================================
\section{Inferência estatística}
\label{sec:met-inferencia}

Os retornos de $h$ dias de datas vizinhas compartilham $h-1$ retornos diários,
e os alvos de 30 dias compartilham até 29. A sobreposição permanece nos dados e
induz autocorrelação nos erros, que não reflete necessariamente dependência
econômica. O erro-padrão HAC de Newey--West não elimina esse problema: ele o
contorna, sob certas condições --- estacionariedade e dependência que decai com
a defasagem ---, tornando a inferência assintoticamente válida. A dissertação
usa o núcleo de Bartlett com $h+1$ defasagens, sendo $h$ o horizonte de
previsão, com pesos $1 - j/(h+2)$, $j = 0, \ldots, h+1$
\citep{newey1987simple}: $h$ defasagens cobrem a sobreposição, e a defasagem
adicional capta a autocorrelação remanescente. Não se aplica correção de
amostra pequena, e o valor-$p$ vem da distribuição normal. A correção atua no
erro-padrão do estimador de MQO, não no estimador. Mesmo com ela, 987 datas com
alvos de 30 dias sobrepostos equivalem a cerca de 33 observações
independentes, e os testes têm pouco poder.

Nos Capítulos~\ref{chap:retornos_linear} e~\ref{chap:retornos_nao_linear}, o regressor é
gerado: o BVRP previsto vem de um primeiro estágio estimado, e o erro-padrão de
MQO do segundo estágio ignora a incerteza dessa estimação \citep{pagan1984}.
Por isso, a inferência principal usa um bootstrap em blocos em dois níveis
\citep{kunsch1989}. No primeiro nível, o bootstrap reestima o modelo de
previsão, em cada origem, numa amostra reamostrada em blocos dentro da própria
janela de treino, o que preserva o embargo. No segundo, reamostra em blocos os
pares (previsão, retorno) das 987 datas e reestima $\beta_h$. O bootstrap em
blocos tem validade assintótica sob condições sobre o tamanho do bloco, que
cresce com a amostra, e pode trazer refinamentos em amostras finitas em alguns
casos; ele não dispensa hipóteses assintóticas. O
Capítulo~\ref{chap:retornos_linear} detalha o procedimento. Como a dissertação
testa cada especificação em seis horizontes, ela reporta também quais
coeficientes sobrevivem à correção de Bonferroni ($0{,}05/6 \approx 0{,}0083$).
Como os seis horizontes são fortemente correlacionados entre si, a dissertação
também testa conjuntamente a hipótese de que $\beta_h = 0$ em todos os
horizontes, com a matriz de covariância dos coeficientes obtida no bootstrap;
os Capítulos~\ref{chap:retornos_linear} e~\ref{chap:retornos_nao_linear} detalham o teste.

% =========================================================
\section{Avaliação fora da amostra}
\label{sec:met-avaliacao}

O Capítulo~\ref{chap:previsao_bvrp} mede o desempenho das previsões pelo $R^2$
fora da amostra contra a média histórica \citep{campbell2008predicting},
\[
R^2_{\text{fora}} = 1 - \frac{\sum_t (y_t - \hat{y}_t)^2}{\sum_t (y_t - \bar{y}_t)^2},
\]
em que $\hat{y}_t$ é a previsão do modelo e $\bar{y}_t$ é a média dos alvos já
conhecidos na origem vigente ($s \le t - h$). Um $R^2_{\text{fora}}$ positivo
indica que o modelo erra menos que a média histórica. Dois procedimentos
complementam a métrica. O teste de \citet{clark2007approximately} compara cada
modelo com a média histórica, que é um modelo aninhado, contra a alternativa
unilateral de que o modelo é melhor. Ele regride uma série de diferenças de
perda numa constante, com erro-padrão HAC de $h+1$ defasagens
(Seção~\ref{sec:met-inferencia}).

O \textit{model confidence set} (MCS) de \citet{hansen2011model} compara todos
os modelos ao mesmo tempo. A partir do conjunto inicial, o procedimento testa a
hipótese de que todos têm a mesma perda esperada. Se a hipótese é rejeitada,
elimina o modelo de pior desempenho relativo e repete o teste com os que
restam. O conjunto final contém, assintoticamente e com probabilidade de pelo
menos $1-\alpha$, os modelos de menor perda esperada. Cada modelo recebe um
valor-$p$ e pertence ao conjunto de nível $\alpha$ sempre que esse valor-$p$ é
maior ou igual a $\alpha$. A dissertação usa a perda quadrática e a estatística
$T_{\max}$, o maior desvio padronizado entre a perda média de um modelo e a do
conjunto. A distribuição vem de um bootstrap com os mesmos blocos móveis de 60
dias do segundo nível do bootstrap da Seção~\ref{sec:met-inferencia}, mais
longos que a sobreposição de 30 dias dos alvos, com 9.999 reamostragens; o
nível principal é de 10\%. Como sensibilidade, a dissertação repete o
procedimento com blocos de 30 e 90 dias, com a estatística $T_R$ (o maior
desvio padronizado entre dois modelos) e com o nível de 25\%. O conjunto inicial reúne os cinco modelos e três
modelos de referência: a média histórica, a proxy e o prêmio realizado
defasado. O teste de \citet{diebold1995comparing}, com perda quadrática,
alternativa bilateral e o mesmo erro-padrão HAC, fica para comparações entre
dois modelos específicos.

% =========================================================
\section{Regimes de volatilidade}
\label{sec:met-regimes}

O regime de alta volatilidade é definido por $D_t = 1$ se
$\text{VH}_{t-29:t} > 37{,}3\%$ a.a. O corte vem da mistura de duas normais
estimada apenas na primeira janela de estimação (Seção~\ref{sec:dados-regimes}), de
modo que nenhuma data é classificada com informação posterior a ela. Como
robustez, a dissertação reestima o corte em cada origem, só com os dados
disponíveis até ela (janela expansiva). No modelo em árvore do
Capítulo~\ref{chap:retornos_nao_linear}, a própria floresta escolhe os cortes, que o capítulo
compara com o de $37{,}3\%$ a.a.
